\section{A supported Gaussian integral}

\begin{theorem}[Supported Gaussian approximation]
    \label{thm:area_supported_gaussian_integral}
    \lean{TNLean.PEPS.AreaLaw.exists_supported_gaussian_integral,
        TNLean.PEPS.AreaLaw.supported_gaussian_integral}
    \leanok
    \uses{def:area_law_hamiltonian,def:qc_site_expectation,
        def:area_graph_ball_localization_constants}
    To establish the supported Gaussian approximation used in
    \cite{OpenAI2026PolynomialPEPS}, fix integers $q\geq1$ and $R\geq0$,
    and a real number $J$.
    There are positive constants $C,v,c$ such that the following holds.
    Let $\Lambda\subseteq\Z^2$ be a finite domain, let $H'$ be a
    Hamiltonian of range $R$ and term-norm bound $J$ on $\Lambda$, and
    let $H$ be any Hermitian operator on the same space.
    Suppose that $W$ acts on $Q\subseteq\Lambda$ and that $D=H'-H$
    acts on $Z\subseteq\Lambda$. For each integer $l\geq0$, put
    $B=N_l(Q\cup Z)$, where distance is measured in the induced graph.
    There is a differentiable path $U$ such that
    \begin{align*}
        U(0)  & =1,                            \\
        U'(t) & =iE_B(e^{itH'}De^{-itH'})U(t).
    \end{align*}
    For every real $t$, the operator $U(t)$ is unitary, acts on $B$,
    and is chosen before the Gaussian variance and time cutoff.
    For every $\eta\geq0$ and $T\geq0$, let $\gamma_\eta$ be the centered
    Gaussian measure of variance $\eta$, and define
    \begin{align}
        M_{\eta,T}
         & =\int_{[-T,T]}E_B(e^{itH'}We^{-itH'})U(t)\,d\gamma_\eta(t),
        \label{eq:supported_gaussian_integral}
        \\
        \widetilde M_{\eta,T}
         & =\int_{[-T,T]}e^{itH'}We^{-itH}\,d\gamma_\eta(t).
        \label{eq:supported_gaussian_original}
    \end{align}
    Then $M_{\eta,T}$ acts on $B$, its norm is at most $\|W\|$, and
    \begin{align*}
        \|M_{\eta,T}-\widetilde M_{\eta,T}\|
        \leq \|W\|C|\Lambda|(1+\|D\|T)e^{vT-cl}.
    \end{align*}
    The integrals retain their original Gaussian mass, and zero variance,
    empty domains, and empty regions are included. If $J\geq0$, the
    estimate holds with $C=2C_0(R)$, $v=v(R,J)$ and $c=1/2$, the constants
    of Definition~\ref{def:area_graph_ball_localization_constants}.
\end{theorem}
\begin{proof}\leanok
    \uses{thm:area_region_localization,
        thm:qc_supported_projected_relative_evolution,
        lem:qc_site_expectation_properties,
        def:area_graph_ball_localization_constants}
    Take constants $C,v,c$ of Theorem~\ref{thm:area_region_localization}
    with $C,v\geq0$; for $J\geq0$ these may be $2C_0(R)$, $v(R,J)$ and $1/2$.
    Applying
    Theorem~\ref{thm:qc_supported_projected_relative_evolution} to the
    projection region $B$ constructs the path $U$ from $H',H,B$
    independently of $\eta$ and $T$. Setting
    $A_B(t)=E_B(e^{itH'}We^{-itH'})$, both $A_B(t)$ and $U(t)$ act on
    $B$, and expectation contractivity together with unitarity gives
    $\|A_B(t)U(t)\|\leq\|W\|$.
    Continuity and this bound imply Gaussian integrability.
    The expectation onto $B$ commutes with integration and fixes the
    integrand in (\ref{eq:supported_gaussian_integral}), proving the
    asserted support, while the restricted Gaussian mass is at most one,
    which proves the norm bound.

    Both $W$ and $D$ act on $Q\cup Z$. For $|t|\leq T$, the region
    estimate and the relative-evolution comparison give
    \begin{align*}
        \|e^{itH'}We^{-itH'}-A_B(t)\|
         & \leq \|W\|C|\Lambda|e^{vT-cl},   \\
        \|e^{itH'}e^{-itH}-U(t)\|
         & \leq \|D\|C|\Lambda|T e^{vT-cl}.
    \end{align*}
    Factoring $e^{itH'}We^{-itH}$ as
    $(e^{itH'}We^{-itH'})(e^{itH'}e^{-itH})$ and subtracting
    $A_B(t)U(t)$ in this order yields the claimed error for the
    integrands by submultiplicativity. Integrating and using again that
    the restricted mass is at most one proves the comparison with
    (\ref{eq:supported_gaussian_original}).
\end{proof}
